\begin{lemma}[7.1.7]
\label{IV.2.s7.1.7-fr}
Sean $A$ un anillo local noetheriano y $\mathfrak p_i$
($1\leqslant i\leqslant r$) sus ideales primos minimales.
Supongamos que cada uno de los anillos
$A/\mathfrak p_i$ es formalmente equidimensional.
Entonces, para todo ideal primo $\mathfrak q$ de $A$
y todo $i$ tal que $\mathfrak p_i\subset\mathfrak q$,
el anillo $A_{\mathfrak q}/\mathfrak p_iA_{\mathfrak q}$
es formalmente equidimensional.
\end{lemma}

\begin{proof}
Como $A_{\mathfrak q}/\mathfrak p_iA_{\mathfrak q}$ es
el anillo local de $A/\mathfrak p_i$ en el ideal primo
$\mathfrak q/\mathfrak p_i$, es posible limitarse al
caso en que $A$ es íntegro y formalmente
equidimensional. Pongamos $A'=\widehat A$ y sea
$\mathfrak q'$ uno de los ideales primos de $A'$
minimales entre los que contienen a $\mathfrak qA'$.
Si se pone $B=A_{\mathfrak q}$,
$B'=A'_{\mathfrak q'}$ es un $B$-módulo plano
(\textbf 0$_{\mathrm I}$, 6.3.2). Pongamos $C=\widehat B$
y $C'=\widehat{B}'$. Como $B'$ es un $B$-módulo
plano, se sabe que $C'$ es un $C$-módulo plano
(Bourbaki, \emph{Alg. comm.}, cap. III, § 5,
n.\textsuperscript{o} 4, prop. 4). Como $A'$ es
catenario (5.6.4) y equidimensional por hipótesis,
también lo es $B'=A'_{\mathfrak q'}$ (\textbf 0, 16.1.4).
Además, por el teorema de Cohen (\textbf{0}, 19.8.8),
$A'$ es isomorfo a un cociente de un anillo regular,
y lo mismo vale para $B'$ (\textbf 0, 17.3.9). Se deduce,
pues, de (7.1.5) que $C'$ es equidimensional.
Por otra parte, $C'$ es catenario (5.6.4); por tanto,
$C$ es equidimensional en virtud de (7.1.3, (i)).
\end{proof}
